\chapter{Tâche 1 — Outil de duplication de configuration Grafana}
\label{chap:t1}
% Entreprise : Free Mobile

% =====================================================================
%  Trame conseillée (réutilisée pour chaque tâche) :
%  1. Contexte & dette technique de départ
%  2. Objectif de la mission
%  3. Conception / choix techniques
%  4. Réalisation (extraits de code COMMENTÉS, schémas — PAS de code entier)
%  5. Difficultés rencontrées
%  6. Résultats & bénéfice business (réduction de la dette)
% =====================================================================

\section{Contexte et dette technique initiale}
La supervision de la valorisation s'appuie sur \textbf{Grafana}, qui agrège dans
des \emph{dashboards} les métriques remontées par le \emph{Valo Engine}. Or cette
supervision n'existe pas sur une seule instance : on en retrouve plusieurs
— typiquement une par environnement (développement, recette, production) ou par
réseau — qui doivent présenter les \textbf{mêmes tableaux de bord}. Faire vivre
ces instances en parallèle suppose donc de \textbf{répliquer} la configuration de
l'une vers les autres à chaque évolution.

Cette réplication ne se résume pas à une copie : chaque tableau de bord référence
ses \textbf{sources de données} (\emph{datasources}) par un \textbf{identifiant
propre à l'instance}. Le même tableau, copié tel quel d'une instance à l'autre,
pointe alors vers des identifiants qui n'y existent pas, et l'affichage se
retrouve cassé. La synchronisation manuelle — export, copie, puis correction au
cas par cas des identifiants — était donc à la fois \textbf{fastidieuse et source
d'erreurs} : une dette d'outillage sur la chaîne de supervision elle-même.

\section{Objectif de la mission}
La mission consistait à \textbf{reprendre le travail d'un stagiaire}, qui avait
amorcé un \gls{poc} sur ce sujet, pour en faire un outil exploitable. Le but :
\textbf{dupliquer l'ensemble des tableaux de bord} d'une instance Grafana vers une
autre, puis \textbf{réécrire automatiquement} les identifiants de sources de
données pour qu'ils correspondent à ceux de l'instance cible. Le cœur du problème
est donc un \textbf{appariement} : faire correspondre, deux à deux, les sources
de données de l'instance d'origine et celles de l'instance de destination, afin
de substituer chaque identifiant par son équivalent.

\section{Conception et choix techniques}
La logique métier est directe : Grafana expose une \gls{api} HTTP permettant de
\textbf{lister} les tableaux de bord et les sources de données d'une instance,
puis d'y \textbf{créer} de nouveaux objets. L'outil interroge l'instance source,
construit une \textbf{table de correspondance} entre les sources de données des
deux instances (appariées par nom), puis parcourt chaque tableau de bord pour y
\textbf{remplacer} les identifiants avant de le pousser sur l'instance cible.

L'enjeu de conception portait surtout sur l'\textbf{interface}. Plutôt qu'une
ligne de commande aride, l'outil a été doté d'une \textbf{interface en mode texte}
(\gls{tui}) à l'aide du framework Python \textbf{Textual}. Ce framework a été une
\textbf{découverte marquante} : il permet de décrire une interface terminal selon
un modèle très proche du développement web — une structure de composants
apparentée au \emph{HTML}, mise en forme par des feuilles de style apparentées au
\emph{CSS}. On obtient ainsi, dans un simple terminal, une interface interactive
(listes, sélection des instances, validation de l'appariement) sans quitter
l'environnement de travail des administrateurs.

\section{Réalisation}
% TODO : remplacer par un extrait réel du dépôt (illustratif, pas le code entier).
Le cœur de l'outil tient dans la réécriture des identifiants : à partir de la
table de correspondance construite entre les deux instances, chaque référence à
une source de données présente dans un tableau de bord est substituée par son
équivalent côté cible.

\begin{lstlisting}[language=Python,caption={Réécriture des identifiants de sources de données lors de la duplication (extrait illustratif).},label={lst:t1}]
# correspondance : datasource de l'instance source -> datasource de la cible,
# appariées par leur nom (ex. "prometheus-valo" présent sur les deux instances)
correspondance = apparier_par_nom(sources_origine, sources_cible)

def reecrire(dashboard):
    # on substitue chaque identifiant de datasource par son équivalent cible ;
    # un dashboard copié sans cette étape pointerait vers des IDs inexistants
    for panel in dashboard["panels"]:
        ancien = panel["datasource"]["uid"]
        panel["datasource"]["uid"] = correspondance[ancien]
    return dashboard
\end{lstlisting}

\section{Difficultés rencontrées}
S'agissant de ma \textbf{première tâche} dans l'équipe, l'enjeu principal n'était
pas tant la complexité — le besoin était circonscrit et la mission relativement
\textbf{rapide} — que la \textbf{prise en main} de l'environnement : le mode de
travail de l'équipe, le dépôt, et la reprise d'un \gls{poc} déjà entamé par un
tiers. La principale subtilité technique tenait à la \textbf{robustesse de
l'appariement} des sources de données, l'appariement par nom supposant une
nomenclature cohérente d'une instance à l'autre.

\section{Résultats et bénéfice business}
% TODO : quantifier si possible (temps de synchronisation avant/après, erreurs
%        évitées). Relier au thème dette technique (chap.~\ref{chap:veille}).
L'outil a transformé une synchronisation manuelle et faillible en une opération
\textbf{automatisée et reproductible}, supprimant la classe d'erreurs liée aux
identifiants de sources de données mal réécrits. Au-delà du livrable, cette
première mission a joué un rôle de \textbf{mise en jambe} : prise de contact avec
l'équipe et son environnement, et découverte de Textual, un framework réinvesti
par la suite. \emph{[À compléter : gain de temps mesuré sur la synchronisation
des instances.]}
